% ============================================================
%  Candidature EDITE - MORANDEAU Timothée
%  Beamer presentation  •  Metropolis theme
% ============================================================
\documentclass[aspectratio=169, 12pt]{beamer}

% ---------- Theme & colours ---------------------------------
\usetheme{metropolis}
\usepackage{appendixnumberbeamer}
\usepackage{graphicx}
\usepackage{microtype}
\usepackage{lmodern}
\microtypesetup{expansion=false}

\definecolor{LIP6blue}{HTML}{1B3A5C}
\definecolor{LIP6teal}{HTML}{2A9AB5}
\definecolor{LIP6gray}{HTML}{7A8FA6}
\definecolor{AccentOrange}{HTML}{E8652A}
\definecolor{AccentGreen}{HTML}{5DAD6F}
\definecolor{BgLight}{HTML}{F5F7FA}

\setbeamercolor{normal text}{fg=LIP6blue, bg=BgLight}
\setbeamercolor{frametitle}{bg=LIP6blue, fg=white}
\setbeamercolor{progress bar}{fg=LIP6teal, bg=LIP6gray!30}
\setbeamercolor{title separator}{fg=LIP6teal}
\setbeamercolor{alerted text}{fg=LIP6teal}
\setbeamercolor{example text}{fg=AccentGreen}
\setbeamercolor{block title}{bg=LIP6blue!90, fg=white}
\setbeamercolor{block body}{bg=LIP6blue!8, fg=LIP6blue}
\setbeamercolor{block title example}{bg=AccentGreen!80, fg=white}
\setbeamercolor{block body example}{bg=AccentGreen!10, fg=LIP6blue}

\metroset{
  titleformat=smallcaps,
  sectionpage=progressbar,
  subsectionpage=none,
  progressbar=frametitle,
  block=fill
}

% ---------- Packages ----------------------------------------
\usepackage[T1]{fontenc}
\usepackage[utf8]{inputenc}
\usepackage[french]{babel}
\usepackage{microtype}
\usepackage{booktabs}
\usepackage{tikz}
\usetikzlibrary{positioning, arrows.meta, shapes.geometric, calc,
                decorations.pathreplacing, backgrounds, fit}
\usepackage{xcolor}
\usepackage{tcolorbox}
\tcbuselibrary{skins, breakable}

% ---------- Custom commands ---------------------------------
\newcommand{\kswap}{\textit{k-swap}}
\newcommand{\egograph}{\textit{ego-graphe}}
\newcommand{\highlight}[1]{\textbf{\textcolor{LIP6teal}{#1}}}

% Thin coloured rule under frame title
\makeatletter
\setlength{\metropolis@titleseparator@linewidth}{1.5pt}
\makeatother

\setbeamertemplate{bibliography item}{}

% ============================================================
%  TITLE PAGE
% ============================================================
\title{\Large Modèles neutres\\ de réseaux complexes intégrant \\ des contraintes temporelles\\ et d'ordre supérieur}
\subtitle{Candidature EDITE}
\author{\textbf{Timothée Morandeau}}
\date{Mai 2026}
\institute{%
  LIP6 - ComplexNetworks, Sorbonne Université
}
\titlegraphic{
  \begin{flushright}
    \includegraphics[height=1.2cm]{sorbonne-logo.png}\hspace{0.1cm}

    \includegraphics[height=1.2cm]{lip6-logo.png}
  \end{flushright}
}

% ============================================================
\begin{document}
% ============================================================

\maketitle

% ----- Plan -------------------------------------------------
\begin{frame}{Plan}
  \tableofcontents
\end{frame}

% ============================================================
\section{Parcours académique}
% ============================================================

\begin{frame}{Baccalauréat Général \hfill{\small 2018-2021}}
  \begin{columns}[c]
    \begin{column}{0.55\textwidth}
      \textbf{\textit{Lycée Charles-Baudelaire, Fosses}}
      \medskip

      \begin{itemize}
        \item \highlight{Mathématiques}
        \item \highlight{Physique-Chimie}
        \item Sciences de l'ingénieur (Première)
      \end{itemize}
    \end{column}
  \end{columns}
\end{frame}

% -----------------------------------------------------------
\begin{frame}{Licence d'Informatique \hfill{\small 2021-2024}}
  \textit{\textbf{Université Paris-Saclay}}
  \medskip

  \begin{columns}[T]
    \begin{column}{0.48\textwidth}
      \begin{itemize}
        \item Sciences des données
        \item Génie Logiciel
        \item Algorithmique
      \end{itemize}
    \end{column}
    \begin{column}{0.48\textwidth}
      \begin{itemize}
        \item Programmation (impérative, OO, fonctionnelle)
        \item Systèmes d'exploitation
        \item \alert{Option L3 :} Informatique théorique
      \end{itemize}
    \end{column}
  \end{columns}
\end{frame}

% -----------------------------------------------------------
\begin{frame}{Master Sciences et Technologies du Logiciel \hfill{\small 2024-2026}}
  \textit{\textbf{Sorbonne Université - filière Recherche}}
  \medskip

  \begin{columns}[T]
    \begin{column}{0.48\textwidth}
      \begin{block}{Cours}
        \begin{itemize}
          \item Algorithmique avancée
          \item Théorie des Graphes
          \item Langages et Compilation
          \item Concurrence
          \item Validation de Programmes
          \item Génie Logiciel
        \end{itemize}
      \end{block}
    \end{column}
    \begin{column}{0.48\textwidth}
      \begin{exampleblock}{Projets de recherche}
        \begin{itemize}
          \item \textbf{M1} - Génération aléatoire d'entiers\\
            {\small \textit{Encadré par A.~Genitrini \& M.~Naima}}
          \item \textbf{M2} - Théorie des catégories\\
            et ses applications en informatique\\
            {\small \textit{Projet exploratoire}}
        \end{itemize}
      \end{exampleblock}
    \end{column}
  \end{columns}
\end{frame}

% ============================================================
\section{Stage de M2 - \textit{Realistic Models of Personal Networks}}
% ============================================================

\begin{frame}{Contexte du stage \hfill{\small Mai-Août 2026}}
  \begin{columns}[c]
    \begin{column}{0.58\textwidth}
      \begin{block}{Où ?}
        LIP6, Sorbonne Université\\
        Équipe \highlight{ComplexNetworks}
      \end{block}
      \medskip
      \begin{block}{Encadrement}
        Mehdi Naima \& Lionel Tabourier
      \end{block}
      \medskip
      \begin{block}{Projet}
        CNRS Miti \highlight{MonReso}
      \end{block}
    \end{column}
    \begin{column}{0.38\textwidth}
        \begin{figure}
            \centering
            \includegraphics[width=0.4\linewidth]{MITI_sansbg.png}
            \label{fig:placeholder}
        \end{figure}
      \begin{tcolorbox}[
        colback=LIP6teal!12, colframe=LIP6teal,
        title={\textbullet\quad Miti MonReso},
        fonttitle=\small\bfseries,
        rounded corners
      ]
        \small Programme du CNRS dédié à l'analyse \emph{interdisciplinaire} des
        réseaux sociaux, à l'interface entre informatique et sociologie.
      \end{tcolorbox}
    \end{column}
  \end{columns}
\end{frame}

% -----------------------------------------------------------
\begin{frame}{Objectif du stage}
  \begin{alertblock}{Question centrale}
    Construire des \highlight{modèles neutres réalistes} d'ego-graphes en
    déterminant les \emph{motifs structurants} des ces réseaux, en étendant la methode PKS \cite{tabourier2024}
    pour contraindre certains motifs.
  \end{alertblock}

  \bigskip
  \begin{columns}[T]
    \begin{column}{0.46\textwidth}
      \begin{block}{\(\rightarrow\)\quad Modèle neutre}
        Graphes aléatoires respectant un ensemble de propriétés fixées
        (distribution des degrés, motifs\dots)
      \end{block}
    \end{column}
    \begin{column}{0.46\textwidth}
      \begin{block}{\(\rightarrow\)\quad Ego-graphe}
        Graphe centré sur un individu, ses relations directes et les liens
        entre ces relations.
      \end{block}
    \end{column}
  \end{columns}
\end{frame}

\begin{frame}{Exemple d'ego-graphe -- réseau Facebook}

  \begin{columns}[c]

    \begin{column}{0.52\textwidth}
      \centering
      \includegraphics[width=\linewidth]
      {Capture d’écran 2026-04-23 143620(1)(1).png}
    \end{column}

    \begin{column}{0.44\textwidth}

      Graphe des amis Facebook d’un individu et des relations
      entre ces amis.

    \end{column}

  \end{columns}
  \scriptsize\textcolor{LIP6gray}{
        Raphael Charbey and Christophe Prieur. \textit{Stars, holes, or paths across your Facebook friends : A graphlet-based characterization of many networks.} Network Science, 7(4) :476–497, December 2019 
      }

\end{frame}

\begin{frame}{Pourquoi les modèles neutres ?}
      \begin{alertblock}{Utilité}
      \begin{itemize}      
        \item[\(\rightarrow\)] Déterminer quelle part des propriétés observées dans un réseau réel
        est produite par l'aléatoire ou
        \emph{structurante} (révélatrice d'un comportement social)
        \item[\(\rightarrow\)] Des modèles neutres réalistes peuvent permettre des simulations contrôlées
    \end{itemize}
      \end{alertblock}
\end{frame}

% -----------------------------------------------------------


% -----------------------------------------------------------
\begin{frame}{Tâches et compétences mobilisées}

      \begin{block}{Compétences}
        \begin{itemize}
          \item Génération aléatoire de graphes
          \item Sampling MCMC (\kswap{} probabiliste)
          \item Analyse de données de réseaux réels
          \item Comptage de motifs (graphlets)
        \end{itemize}
      \end{block}

\end{frame}

% ============================================================
\section{Intérêts de recherche}
% ============================================================

\begin{frame}{Domaines d'intérêt}

      \begin{itemize}
        \setlength\itemsep{8pt}
        \item \highlight{Algorithmique} - complexité, structures de données
        \item \highlight{Théorie des Graphes}
        \item \highlight{Approche interdisciplinaire} 
        \item \highlight{Analyse de données}
      \end{itemize}

\end{frame}

% -----------------------------------------------------------
\begin{frame}{Structures de données complexes}
  \begin{columns}[T]
    \begin{column}{0.48\textwidth}
      \begin{block}{Hypergraphes}
        Généralisation des graphes où une \emph{hyperarête} peut connecter
        \emph{plus de deux} sommets simultanément.
      \end{block}
      \begin{block}{Projet M2}
        UE AAGA (Algorithmique avancées, Graphes et Applications) :\\
        \textit{Mesures d'homophilie dans les graphes et hypergraphes}
      \end{block}
    \end{column}
    \begin{column}{0.48\textwidth}
        \begin{figure}
            \centering
            \includegraphics[width=0.5\linewidth]{patate(1).png}

            \label{fig:placeholder}
        \end{figure}
    \end{column}
  \end{columns}

\end{frame}


% -----------------------------------------------------------


% ============================================================
\section{Motivations et projet doctoral}
% ============================================================
\begin{frame}{Aspirations professionnelles}
  \begin{itemize}
    \setlength\itemsep{8pt}
    \item Poursuivre en \highlight{recherche académique} après la thèse
    \item Contribuer à la recherche en algorithmique et théorie des graphes
  \end{itemize}
\end{frame}

\begin{frame}{Pourquoi ce cadre ?}
  \begin{columns}[T]
    \begin{column}{0.48\textwidth}
      \begin{block}{Continuité avec le stage}
        \begin{itemize}
          \item Maîtrise du \kswap{} probabiliste et des modèles neutres
                d'ego-graphes
          \item Analyse de motifs (graphlets) sur données réelles
        \end{itemize}
      \end{block}
    \end{column}
    \begin{column}{0.48\textwidth}
      \begin{exampleblock}{Équipe ComplexNetworks}
        \begin{itemize}
          \item L.~Tabourier - co-auteur de la méthode PKS \cite{tabourier2024}
          \item M.~Naima - graphes temporels et centralité \cite{naima2023}
          \item Projet MonReso - environnement interdisciplinaire
        \end{itemize}
      \end{exampleblock}
    \end{column}
  \end{columns}
\end{frame}

%--------------------------------------------------------

\begin{frame}{Problème scientifique}
  \begin{alertblock}{Question centrale}
    Comment générer des modèles \highlight{réalistes et non-biaisés} pour des structures
    plus expressives : graphes \emph{temporels} et \emph{d'ordre supérieur} ?
  \end{alertblock}

  \bigskip

  \begin{block}{Enjeux}
    \begin{itemize}
      \item Comparer PKS et ERGM sur les graphes statiques
      \item Etendre la méthode PKS aux graphes temporels et d'ordre supérieur
      \item Identifier les motifs structurants de ces réseaux complexes
    \end{itemize}
  \end{block}
\end{frame}

% -----------------------------------------------------------
{
\begin{frame}{Conclusion}
  \textcolor{LIP6teal}{\normalsize Le sujet}\\[0.4em]
  \small Des questions ouvertes et importantes en science des réseaux

  \vspace{1.2em}

  \textcolor{LIP6teal}{\normalsize Ma candidature}\\[0.4em]
  \small Un parcours qui s'inscrit naturellement dans ce projet
\end{frame}
}
% ============================================================
%  CLOSING SLIDE
% ============================================================
{
\setbeamercolor{background canvas}{bg=LIP6blue}
\begin{frame}[standout]
  \textcolor{white}{\Large Merci pour votre écoute !}\\[1em]
  \textcolor{LIP6teal}{\normalsize Questions ?}
\end{frame}
}

% ============================================================
%  REFERENCES
% ============================================================
\appendix
\begin{frame}[allowframebreaks]{Références}
  \small
  \begin{thebibliography}{99}

    \bibitem{tabourier2024}
    L.~Tabourier and J.~Karadayi.
    \newblock Probabilistic $k$-swap method for uniform graph generation beyond
      the configuration model.
    \newblock \textit{Journal of Complex Networks}, 12(1), February 2024.

    \bibitem{charbey2019}
    R.~Charbey and C.~Prieur.
    \newblock Stars, holes, or paths across your Facebook friends: A
      graphlet-based characterization of many networks.
    \newblock \textit{Network Science}, 7(4):476-497, December 2019.

    \bibitem{gaumont2018}
    N.~Gaumont, M.~Panahi, and D.~Chavalarias.
    \newblock Reconstruction of the socio-semantic dynamics of political activist
      Twitter networks: method and application to the 2017 French presidential
      election.
    \newblock \textit{ICWSM}, 2018.

    \bibitem{naima2023}
    M.~Naima. 
    \newblock Temporal Betweenness Centrality on Shortest Walks Variants. 
    \newblock \textit{Applied Network Science}, 2025.


  \end{thebibliography}
\end{frame}

\end{document}
